\documentclass[9pt,twoside]{extarticle}
\usepackage{Breviary-Appendix}
\usepackage{indentfirst}
%%
\begin{document}
\thispagestyle{empty}
\renewcommand{\thepage}{\arabic{page}*}
%\null
%\vfill
\section*{\center \textsc{Appendix}\\ \normalfont \normalsize {\textsc{\textred{ad usum}\\ \Large Eremitarum Discalceatorum\\ Dominæ Nostræ de Monte Carmelo}}}
\begin{center} \includegraphics[width=.2\textwidth]{images/LogoforHermits-Transparent} \end{center}
\begin{center}\textcolor{black}{\rule{0.4\linewidth}{0.5pt}} \end{center}
\begin{center} {\large \textsc{Sumpta Ex Breviario Carmelitarum 1866\\ \textred{Pars Autumnalis}}} \end{center}
\begin{center}\textcolor{black}{\rule{\linewidth}{0.5pt}} \end{center}

%\vfill

%\newpage

\begin{multicols}{2}
%\begin{center} \includegraphics[width=1 \textwidth]{images/Sacred_Hearts_rotated} \end{center}
%\subsection*{\center \normalfont \textsc{die 10 octobris}}
%%\subsubsection*{\center \textsc{infra octavam}}
%\subsection*{\center In Octava S. Teresiæ\\ a Jesu Infante}
%\subsubsection*{\center \textred{Simplex}}
%
%\small \textred{Commemoratio tantum fit in I Vesperis et ad Laudes:}
%
%\normalsize
%\textred{Magn. Ant.} Ego Dóminus~* sponsábo te mihi in sempitérnum~: et sponsábo te mihi in justítia et judício, et in misericórdia et in miseratiónibus~: et sponsábo te mihi in fide.
%
%\textred{℣.}~Ego autem in Dómino gaudébo. \textred{℟.}~Et exsultábo in Deo Jesu meo.
%
%%\noindent Orémus. \hfill \textred{Oratio}
%\subsubsection*{\center \textred{Oratio}}
%\lettrine{D}ómine, qui dixísti: Nisi efficiámini sicut párvuli, non intrábitis in regnum cælórum: \textbf{†} da nobis, quǽsumus; ita sanctæ Terésiæ Vírginis in humilitáte et simplicitáte cordis vestígia sectári, \textbf{*} ut prǽmia consequámur ætérna: Qui vivis.
%
%\textred{Ben. Ant.} Réquiem tibi dabit~* Dóminus semper, et implébit splendóribus ánimam tuam; et eris quasi hortus irríguus, et sicut fons aquárum, cujus non defícient aquæ.
%
%\textred{℣.}~Erit tibi Dóminus in lucem sempitérnam. \textred{℟.}~Et Deus tuus in glóriam tuam.

%
%\begin{center}\textcolor{red}{\rule{\linewidth}{0.5pt}} \end{center}
\subsection*{\center \normalfont \textsc{die 15 octobris}}
%\subsubsection*{\center \textsc{infra octavam}}
\subsection*{\center {\small De Octava}\\ \textbf{\textsc{S. Teresiæ a Jesu}}}
%\subsubsection*{\center \textred{Duplex minus.}}
%\begin{center}\textred{Vide festum, supra, 535.}\end{center}
\fancyhead[CO]{\textsc{\textred{Infra Octav. S. Teresiæ a Jesu}}}
\fancyhead[CE]{\textsc{\textred{15 Octobris}}}

\smallskip
\small \textred{Per Octavam S. Teresiæ, fit de ea commemoratio in Laudibus per Antiphonam ad Benedictus, et in Vesperis per Antiphonam de II Vesperis Festi.}

\textred{Dominica vero infra Octavam pro Commemoratione dicuntur Antiphonæ de utrisque Vesperis et Laudibus Festi, et ad responsionem Versuum additur} Allelúja.

\normalsize
\textred{Ad Laudes, Ant.} Dedit ei~* Dóminus sapiéntiam et prudéntiam multam nimis, et latitúdinem cordis, quasi arénam, quæ est in líttore maris.

\textred{℣.}~Adjuvábit eam Deus vultu suo. \textred{℟.}~Deus in médio ejus, non commovébitur.

\subsubsection*{\center \textred{Oratio}}
%\begin{center}\textred{Oratio}\end{center}
\lettrine{E}{x}áudi nos, Deus, salutáris noster~: ut, sicut de beátæ Terésiæ Vírginis tuæ et Matris nostræ festivitáte gaudémus; ita cæléstis ejus doctrínæ pábulo nutriámur, et piæ devotiónis erudiámur afféctu. Per Dóminum.

\textred{Ad Vesperas, Ant.} Sapiéntiam ejus~* enarrábunt gentes, et laudem ejus enuntiábit Ecclésia.

\textred{℣.}~Oblivíscere pópulum tuum et domum patris tui. \textred{℟.}~Et concupíscet Rex decórem tuum.

\smallskip
\textred{In I Vesperis Dominicæ, Ant.} Quæsívi~* in sponsam mihi eam assúmere. Doctrix enim est disciplínæ Dei, et eléctrix óperum illíus.

%\begin{center}\textcolor{red}{\rule{\linewidth}{0.5pt}} \end{center}
\subsection*{\center \normalfont \textsc{die 22 octobris}}
%\subsubsection*{\center \normalfont \textsc{Dies Octava}}
\subsection*{\center \textsc{In Octava S. Teresiæ}}
\subsubsection*{\center \textred{{\small Virginis O.N.}\\Duplex minus I classis}}
%\fancyhead[C]{\textsc{22 Oct. - Dies Octava S. Teresiæ.}}

\smallskip
\small \textred{Antiphonæ et Psalmi ad omnes Horas et Versus Nocturnorum de occurrenti Hebdomadæ die, ut in Psalterio; reliqua ut in Festo, 669, præter ea quæ hic habentur propria.}

\textred{In I Vesperis, fit commemoratio præcedentis:}

\normalsize
\textred{Ant.} Hic vir * despíciens mundum et terréna, triúmphans, divítias cœlo cóndidit ore, manu.
\fancyhead[CO]{\textsc{\textred{Dies Octava S. Teresiæ a Jesu}}}
\fancyhead[CE]{\textsc{\textred{Die 22 Octobris}}}

\textred{℣.}~Amávit eum Dóminus et ornávit eum. \textred{℟.}~Stólam glóriæ índuit eum, allelúja.

\smallskip
Orémus. \hfill \textred{Oratio}
\lettrine{O}mnípotens sempitérne Deus, qui, méritis beáti Hilariónis Confessóris tui, semper es et ubíque mirábilis~: quǽsumus cleméntiam tuam; ut sicut ei eminéntem glóriam contulísti, sic ad consequéndam misericórdiam tuam mereámur ejus précibus adjuvári. Per Christum.

{\small \textred{Deinde commemoratio S. Marci, Ep. et Mart.:}}

\textred{Ant.} Iste sanctus * pro lege Dei sui certávit usque ad mortem, et a verbis impiórum non tímuit~: fundátus enim erat supra firmam petram.

\textred{℣.}~Glória et honóre coronásti eum Dómine. \textred{℟.}~Et constituísti eum super ópera mánuum tuárum.

\smallskip
Orémus. \hfill \textred{Oratio}
\lettrine{D}eus, qui nos beáti Marci Mártyris tui atque Pontíficis ánnua solemnitáte lætíficas~: concéde propítius; ut, cujus natalítia cólimus, de ejúsdem étiam protectióne gaudeámus. Per Dóminum.

\subsubsection*{\center \normalfont \textsc{AD MATUTINUM}}
\small \textred{In I Nocturno, Lectiones de Scriptura occurrenti cum suis Responsoriis de Tempore.}

\normalsize
\subsubsection*{\center \normalfont \small \textsc{IN II NOCTURNO}}
\subsubsection*{\center \normalfont Sermo sancti Máximi Epíscopi\\ \small \textred{Hom. in nat. S. Agnetis\\
\normalsize Lectio iv}}
\lettrineThree{C}{um} in toto mundo virgíneus flos Maríæ immarcescíbiles corónas innéctat, et sceptrígeram pudóris aulam immaculáto consérvet afféctu, eoúsque intégritas perseverávit ad palmam, ut in puéllis trophæum sanctitátis arríperet, et per vestígia matris Vírginis ad cœléstem thálamum perveníret. O quam imitábile virgínibus, Virgo, tui amóris tradidísti exémplum! O quam sanctum pectóribus dandum virgíneis tribuísti respónsum, contemnéndo mundi divítias, exsecrándo carnis concupiscéntias, et solam Christi pulchritúdinem adamándo! Accédite ad puéllam puéllæ, et díscite, qualis in péctore suo circa amórem Christi flammas accéndit. Díscite, vírgines, Christi amórem in ea fervéntem, et omnes mundi divítias velut stércora recusántem.

\textred{℟.}~Dixit mihi Dóminus meus~: Ne tímeas fília~: \textred{*}~Non rápiet te quisquam de manu mea. \textred{℣.}~Neque mors, neque vita, neque creatúra ália póterit me separáre a caritáte Dei. \textred{N}on.

\subsubsection*{\center \normalfont \textred{Lectio v}}
\lettrine{I}deo enim sanctórum per lítteras tradúntur exémpla, ut unusquísque pro qualitáte sexus modum sui propósiti téneat, et sicut per vírginem vírgines, ita per cónjugem cónjuges víncere sæculum pro Christi amóre stúdeant, ac repugnatióne ánimi, mundi dúlcia recusántes, et ejus non recusántes austéra, ad delectatiónes perpétuas, et ad gáudia ætérna pertíngant. Hæc dicéntes, non vituperámus cónjugum casta commércia, sed vírginum perseverántiam adornámus. Non enim virgínitas malis rebus superponénda est, sed est rebus óptimis præferénda~: quæ enim mundis múndior, et sanctis est sánctior, non índiget malárum rerum vituperatióne laudári.

\textred{℟.}~Gaudens gaudébo in Dómino, et exsultábit ánima mea in Deo meo, quia índuit me vestiméntis salútis; \textred{*}~Et induménto justítiæ circúmdedit me. \textred{℣.}~Quasi sponsam decorátam coróna, et quasi sponsam ornátam monílibus suis. \textred{E}t.

\subsubsection*{\center \normalfont \textred{Lectio vj}}
\lettrine{V}eni jam, Virgo, ad thálamum, quem quæsísti, ut sponsam te virgínitas immaculáta cómprobet Christi. Felix catérva vírginum, quibus stúdium est tuæ sequi confessiónis vestígia. Habébunt enim tecum corónam in cœlo, quæ tecum superavérunt inimícum in mundo. Nam sicut tibi vestígia Maríæ prosequénti participátio ætérnæ remuneratiónis appáruit~: ita eis te imitántibus ætérna crédimus gáudia non negánda. Itaque o spléndida Christo Dei Fílio, et ómnibus Angelis, atque Archángelis grata! ut nostri meminísse dignéris, quibus póssumus précibus exorámus, ut nobis ille tríbuat indulgéntiam, qui tibi tuórum ómnium labórum trádidit palmam. 

\textred{℟.}~Pulchra fácie, sed púlchrior fide beáta es Terésia virgo, réspuens mundum, lætáberis cum Angelis~: \textred{*}~Intercéde pro ómnibus nobis. \textred{℣.}~Gaudeámus, et exsultémus, et demus glóriam ei, quia venérunt núptiæ Agni. \textred{I}ntercéde. Glória Patri. \textred{I}ntercéde.

\subsubsection*{\center \normalfont \small \textsc{IN III NOCTURNO}}
\subsubsection*{\center \normalfont Lectio sancti Evangélii\\ secúndum Matthǽum\\
\textred{Lectio vij \hfill Cap.~11,~25-30}}
\lettrine{I}n illo témpore~: Respóndens Jesus dixit~: Confíteor tibi Pater, Dómine cœli et terræ, quia abscondísti hæc a sapiéntibus, et prudéntibus, et revelásti ea párvulis. Et réliqua.

\subsubsection*{\center \normalfont Homilía sancti Augustíni Epíscopi\\
\small \textred{Serm. 9. de Verbis Dom.}}

\lettrineThree{N}{on} immérito ille vas electiónis cum ingénti lætítia dixit~: Non sunt condígnæ passíones hujus témporis ad futúram glóriam, quæ revelábitur in nobis. Ecce unde illud jugum suáve est, et sárcina levis, et si angústa est paucis eligéntibus, fácilis tamen ómnibus diligéntibus. Dixit Psalmísta~: Propter verba labiórum tuórum ego custodívi vias duras. Sed quæ dura sunt laborántibus, éadem ipsa mitéscunt amántibus.

\textred{℟.}~Habitáre fecit eam Deus stérilem in domo~: \textred{*}~Matrem filiórum lætántem. \textred{℣.}~Cibávit eam pane vitæ, et intelléctus, et aqua sapiéntiæ salutáris potávit illam. \textred{M}atrem.

\subsubsection*{\center \normalfont \textred{Lectio viij}}
\lettrine{P}ropter quod ita divínæ pietátis dispensatióne actum est, ut intérior homo, qui renovátur de die in diem, non adhuc sub lege pósitus, sed jam sub grátia exonerátus sárcinis innumerabílium observatiónum, quod erat revéra grave jugum, sed duræ cervíci conveniénter impósitum, facilitáte símplicis fídei, et bonæ spei, et sanctæ caritátis, quidquid molestiárum exterióri hómini forínsecus intulísset ille princeps, qui missus est foras, interióri gáudio fíeret leve. Nihil enim tam fácile bonæ voluntáti est, quam ipsa sibi~: et hæc súfficit Deo.

\textred{℟.}~Tu glória Jerúsalem, tu lætítia Israel, tu honorificéntia pópuli nostri, quia fecísti viríliter, et confortátum est cor tuum~: ídeo et manus Dómini confortábit te~: \textred{*}~Et eris benedícta in ætérnum. \textred{℣.}~Surrexérunt fílii tui, et. beatíssimam prædicavérunt. \textred{E}t.

\subsubsection*{\center \normalfont Pro sancto Marco Epíscopo et Mártyre\\
\small \textred{Ex Beda \hfill \normalsize Lectio ix}}
\lettrine{M}arcus Epíscopus, primus ex Géntibus Jerosolymórum Ecclésiam regéndam suscépit. Factum est enim, ut cívitas post intéritum Judáicæ genis, íncolis in eam peregrínæ natiónis confluéntibus, permutátis cívibus, et in Romanórum jus tota convérsa, nomen páriter mutáret, ac ritum. Sacerdótium ígitur Marcus illíus civitátis sanctæ suscípiens, non multo post pro confessióne Christi constantíssime martýrio consummátus, voctóriæ palmam proméruit, undécimo Kaléndas Novémbris.

\textred{℟.}~Super salútem et spéciem diléxi sapiéntiam, et propósui pro luce habére illam~: \textred{*}~Quóniam inexstinguíbile est lumen illíus. \textred{℣.}~Lætáta sum in ómnibus, quóniam antecedébat me ista sapiéntia. \textred{Q}uóniam. Glória Patri. \textred{Q}uóniam.

Te Deum laudámus. %\textred{6}

\textred{Sacerd.~℣.}~Veni, Sponsa Christi, áccipe corónam. \textred{℟.}~Quam tibi Dóminus præparávit in ætérnum, allelúja.

{\small \textred{Ad Laudes, commemoratio S. Marci:}}

\textred{Ant.} Nisi granum * fruménti cadens in terram, mórtuum fúerit~: ipsum solum manet.

\textred{℣.}~Magna est glória ejus in salutári tuo. \textred{℟.}~Glóriam et magnum decórem impónes super eum.
 
\smallskip
Orémus. \hfill \textred{Oratio}
\lettrine{D}eus, qui nos beáti Marci Mártyris tui atque Pontíficis ánnua solemnitáte lætíficas~: concéde propítius; ut, cujus natalítia cólimus, de ejúsdem étiam protectióne gaudeámus. Per Dóminum.

\small \textred{Vesperæ de sequenti, 708, Commemoratio præcedentis:}

\normalsize
\textred{Ant.} Sapiéntiam ejus~* enarrábunt gentes, et laudem ejus enuntiábit Ecclésia.

\textred{℣.}~Oblivíscere pópulum tuum et domum patris tui. \textred{℟.}~Et concupíscet Rex decórem tuum, allelúja.

\subsubsection*{\center \normalfont \textred{Oratio}}
\lettrine{E}{x}áudi nos, Deus, salutáris noster~: ut, sicut de beátæ Terésiæ Vírginis tuæ et Matris nostræ festivitáte gaudémus; ita cæléstis ejus doctrínæ pábulo nutriámur, et piæ devotiónis erudiámur afféctu. Per Dóminum.

\begin{center}\textcolor{red}{\rule{\linewidth}{0.5pt}} \end{center}

\subsection*{\center \normalfont \textsc{die 24 novembris}}
\subsection*{\center {\small De Octava}\\ \textbf{\textsc{S. Joannis a Cruce}}}
%\subsubsection*{\center \textred{Duplex minus.}}
%\begin{center}\textred{Vide festum, supra, 535.}\end{center}
%\fancyhead[C]{\textsc{24 Nov. - Infra Octav. S. Joannis a Cruce}}
\fancyhead[CO]{\textsc{\textred{Infra Octav. S. Joannis a Cruce}}}
\fancyhead[CE]{\textsc{\textred{24 Novembris}}}

\small \textred{Per Octavam S. Joannis, fit de ea commemoratio in Laudibus per Antiphonam ad Benedictus, et in Vesperis per Antiphonam de II Vesperis Festi.}

\textred{Dominica vero infra Octavam pro Commemoratione dicuntur Antiphonæ de utrisque Vesperis et Laudibus Festi, et ad responsionem Versuum additur} Allelúja.

\normalsize
\textred{Ad Laudes, Ant.} Dedit illi * sciéntiam Sanctórum, honestávit illum in labóribus, et complévit labóres illíus.

\textred{℣.}~Justus germinábit sicut lílium. \textred{℟.}~Et florébit in ætérnum ante Dóminum.

\subsubsection*{\center \textred{Oratio}}
\lettrine{D}eus, qui sanctum Joánnem Confessórem tuum, atque Doctórem, Patrem nostrum perféctæ sui abnegatiónis, et crucis amatórem exímium effecísti~: concéde; ut ejus imitatióni júgiter inhæréntes, glóriam assequámur ætérnam. Per Dóminum.

\textred{Ad Vesperas, Ant.} Quæsívit * verba utília, et conscrípsit sermónes rectíssimos, et veritáte plenos.

\textred{℣.}~Ora pro nobis, sancte pater Joánnes. \textred{℟.}~Ut digni efficiámur promissiónibus Christi.

\smallskip
\textred{In I Vesperis Dominicæ, Ant.} Collaudábunt multi * sapiéntiam ejus, et usque in sǽculum non delébitur.

%\begin{center}\textcolor{red}{\rule{\linewidth}{0.5pt}} \end{center}

%\subsection*{\center \textsc{\textred{die 29 novembris}}}
%\subsubsection*{\center \textsc{\textred{Quinta die infra Octav.}}}
%\subsection*{\center \textsc{S. Joannis a Cruce}}
%\subsubsection*{\center \textred{Semiduplex}}
%%\begin{center}\textred{Vide festum, supra, 535.}\end{center}
%\fancyhead[C]{\textsc{29 Nov. - Quinta Die infra Octav. S. Joannis.}}
%\fancyhead[CO]{\textsc{\textred{Quinta Die infra Octav. S. Joannis}}}
%\fancyhead[CE]{\textsc{\textred{29 Novembris}}}
%
%\subsubsection*{\center \textsc{AD MATUTINUM}}
%\small \indent \textred{Invitat. Hymnus, prima Ant. j. Noct. et Psalmi ut in die, 884.}
%
%\textred{Prima et secunda Lectio erit de Homilia in Evang. Vigiliæ S. Andreæ Apostoli, ut infra; tertia lectio de Octava.}
%
%\textred{℣. et Responsorii ut in die, secundum ordinem dierum et Nocturnorum.} \normalsize
%
%\subsubsection*{\textred{Feriis II et V:}}
%\indent \textred{℣.}~Amávit eum Dóminus, et ornávit eum. \textred{℟.}~Stolam glóriæ índuit eum.
%
%\textred{℟. j.} Euge serve bone et fidélis, quia super pauca fuísti fidélis, supra multa te constítuam: \textred{*}~Intra in gáudium Dómini tui. \textred{℣.}~Dómine quinque talénta tradidísti mihi, ecce ália quinque superlucrátus sum. \textred{I}ntra.
%
%\textred{℟. ij.} Justus germinábit sicut lílium: \textred{*}~Et florébit in ætérnum ante Dóminum. \textred{℣.}~Plantátus in domo Dómini, in átriis domus Dei nostri. \textred{E}t.
%
%\textred{℟. iij.} Iste cognóvit justítiam, et vidit mirabília magna, et exorávit Altíssimum. \textred{*}~Et invéntus est in número sanctórum. \textred{℣.}~Iste est, qui contémpsit vitam mundi, et pervénit ad cœléstia regna. \textred{E}t. Glória Patri. \textred{E}t.
%
%\subsubsection*{\textred{Feriis III et VI:}}
%\indent \textred{℣.}~Justum dedúxit Dóminus per vias rectas. \textred{℟.}~Et osténdit illi regnum Dei. 
%
%\textred{℟. j.} Desidérium ánimæ ejus tribuísti ei Dómine. \textred{*}~Et voluntáte labiórum ejus non fraudásti eum. \textred{℣.}~Posuísti Dómine super caput ejus corónam de lápide pretióso. \textred{E}t.
%
%\textred{℟. ij.} Amávit eum Dóminus, et ornávit eum: stolam glóriæ índuit eum: \textred{*}~Et ad portas paradísi coronávit eum. \textred{℣.}~Induit eum Dóminus lorícam fídei, et ornávit eum. \textred{E}t.
%
%\textred{℟. iij.} Iste sanctus digne in memóriam vértitur hóminum: qui ad gáudium tránsiit Angelórum: quóniam in hac peregrinatióne solo córpore constitútus, cogitatióne, et aviditáte, \textred{*}~In illa ætérna pátria conversátus est. \textred{℣.}~Vínculis carnis absolútus, taléntum sibi créditum Dómino suo duplicátum reportávit. \textred{I}n. Glória Patri. \textred{I}n.
%
%\subsubsection*{\textred{Feria IV et Sabbato:}}
%\indent \textred{℣.}~Justus ut palma florébit in domo Dómini. \textred{℟.}~Sicut cedrus Líbani multiplicábitur.
%
%\textred{℟. j.} Magnificávit eum Dóminus in conspéctu regum: \textred{*}~Et dedit illi corónam glóriæ. \textred{℣.}~Státuit illi testaméntum sempitérnum. \textred{E}t.
%
%\textred{℟. ij.} Vir Israelíta, gaude cohéres Christi: \textred{*}~Ora pro nobis. \textred{℣.}~Ut précibus tuis a Deo véniam sperémus, devóti ad te confúgimus. \textred{O}ra.
%
%\textred{℟. iij.} Sancte Joánnes, Christi Conféssor, audi rogántes sérvulos, \textred{*}~Et impetrátam nobis cœ́litus tu defer indulgéntiam. \textred{℣.}~O sancte Joánnes, sidus áureum Dómini grátia, servórum gémitus sólita súscipe cleméntia. \textred{E}t. Glória Patri. \textred{E}t. 
%
%Te Deum laudámus. 6.
%
%\textred{Sacerd. ℣.}~Os justi meditábitur sapiéntiam. \textred{℟.}~Et lingua ejus loquétur judícium.
%
%\subsubsection*{\center \textred{Pro Vigilia S. Andreæ Ap.}}
%\subsubsection*{\center \normalfont Léctio sancti Evangélii secúndum Joánnem.\\
%\textred{\textbf{Lectio j} \hfill Cap. I}}
%\lettrine{I}n illo témpore: Stabat Joánnes, et ex discípulis ejus duo. Et respíciens Jesum ambulántem, dicit: Ecce Agnus Dei. Et réliqua.
%
%\subsubsection*{\center \normalfont Homilía sancti Augustíni Epíscopi.\\
%\small \textred{Tract. 7. in Joann.}}
%\lettrineThree{Q}{u}ia talis erat Joánnes amícus sponsi, non quærébat glóriam suam, sed testimónium perhibébat veritáti. Numquid vóluit apud se remanére discípulos suos, ut non sequeréntur Dóminum? Magis ipse osténdit discípulis suis, quem sequeréntur: habébant enim illum tamquam Agnum. Et ille: Quid me atténditis? Ego non sum Agnus: Ecce Agnus Dei. De quo et supérius díxerat: Ecce Agnus Dei. Et quid nobis prodest Agnus Dei? Ecce, ait, qui tollit peccátum mundi. Secúti sunt illum, noc audíto, duo, qui erant cum Joánne.
%
%\subsubsection*{\center \textred{Lectio ij}}
%\lettrine{V}ideámus sequéntia. Ecce Agnus Dei. Hoc Joánnes. Et audiérunt eum duo discípuli loquéntem, et secúti sunt Jesum. Non sic illum sequebántur, quasi jam ut inhærérent illi: nam maniféstum est, quando illi inhæsérunt, quia de navi eos vocávit. In his enim duóbus erat Andréas, sicut modo audístis. Andréas autem frater Petri erat. Et nóvimus in Evangélio, quod Petrum et Andréam Dóminus de navi vocávit, dicens: Veníte post me, et fáciam vos piscatóres hóminum. Et ex illo jam inhæsérunt illi, ut non recéderent.
%
%\subsubsection*{\center \normalfont Sermo sancti Bernárdi Abbátis\\
%\small \textred{In festo S. Malachiæ. Serm. 2. \hfill \textbf{Lectio~iij}}}
%\lettrine{Q}uasi unus ómnium parens vivébat ómnibus Joánnes, quasi gallína pullos suos, sic fovébat omnes, et in velaménto alárum suárum protegébat. Non sexus, non ætas, non condítio discernebátur, aut persóna: déerat némini, expánso ómnibus grémio pietátis. De quacúmque tribulatióne clamarétur ad eum, própriam reputábat, nisi quod in sua pátiens, in aliéna compátiens erat, plerúmque et impátiens. Non tamen ita ómnibus intendébat, ut se solum expóneret, solum cura excíperet generáli. Erat et sui sollícitus, se ipsum custodiébat. Ita benígne totus suus, et totus ómnium erat, ut nec cáritas a custódia sui, nec propríetas ab utilitáte commúni eum impedíre vel retardáre in áliquo viderétur.
%
%\textred{Ad Laudes, prima Ant. Laudum:} Addúcam eum * in montem sanctum meum, et lætificábo eum in domo oratiónis meæ. 
%
%\textred{Psalmi de Dominica.}
%
%\textred{Cetera ut in die, 888.}
%
%\textred{Et fit com. Vigiliæ per Ant. et ℣. Feriæ occurrentis, ut in Psalterio, cum oratione sequenti:}
%
%\noindent Orémus. \hfill \textred{Oratio}
%\lettrine{Q}uæsumus, omnípotens Deus, ut beátus Andréas Apóstolus, cujus prævenímus festivitátem, tuum pro nobis implóret auxílium, ut a nostris reátibus absolúti, a cunctis étiam perículis eruámur. Per Christum.
%
%\textred{Deinde com. S. Saturnini Mart., Ant.} Nisi granum * fruménti cadens in terram mórtuum fúerit, ipsum solum manet.
%
%\textred{℣.}~Magna est glória ejus in salutári tuo. \textred{℟.}~Glóriam et magnum decórem impónes super eum.
%
%\textred{Oratio} Deus, qui nos beáti Saturníni \textred{ut supra.}
%\item \textred{Vesp. de S. Andrea Apost.} 901.
%\end{tightlist}
%
%\begin{center}\textcolor{red}{\rule{\linewidth}{0.5pt}} \end{center}

\subsection*{\center \normalfont \textsc{die 1 decembris}}}
\subsection*{\center \textsc{In Octava\\ S. Joannis a Cruce}}}
\subsubsection*{\center \textred{{\small Confessoris et Doctoris O.N.}\\ Duplex minus I classis}}
%\begin{center}\textred{Vide festum, supra, 535.}\end{center}
%\fancyhead[C]{\textsc{1 Dec. - Dies Octava S. Joannis a Cruce.}}

\smallskip
\indent \small \textred{Antiphonæ et Psalmi ad omnes Horas et Versus Nocturnorum de occurrenti Hebdomadæ die, ut in Psalterio; reliqua  ut in Festo, 884, exceptis Lectionibus, ut infra.}

\textred{In I Nocturno, Lectiones de Scriptura occurrenti cum suis Responsoriis de Tempore.}

\fancyhead[CO]{\textsc{\textred{Dies Octava S. Joannis a Cruce}}}
\fancyhead[CE]{\textsc{\textred{Die 1 Decembris}}}

\normalsize
\subsubsection*{\center \small \textsc{IN II NOCTURNO}}
\subsubsection*{\center \normalfont Sermo sancti Bernárdi Abbátis\\
\small \textred{In nat. S. Vict. Conf. Serm. 1\\
\normalsize Lectio iv}}
\lettrineThree{J}{o}ánnis vita, et glória speciális, non tam ad glóriam, quam ad virtútem próvocat omnes, qui recto sunt corde. Non recti plane, sed pervérsi est ánimi, ante quǽrere glóriam, quam exercére virtútem, et velle coronári, qui legítime non certáverit. Vanum est vobis, inquit, ante lucem súrgere. Ita est. Frustra ad celsitúdinem nítitur glóriæ qui prius non cláruit virtúte. Absit mihi gloriári, nisi in glória illórum, quibus Prophéta congrátulans, ait~: Dómine, in lúmine vultus tui ambulábunt, et in nómine tuo exsultábunt tota die, et in justítia tua exaltabúntur. Et infert~: Quóniam glória virtútis eórum tu es. Pulchre non eórum, sed virtútis eórum glória commendátur.

\textred{℟.}~Desidérium ánimæ ejus tribuísti ei Dómine. \textred{*}~Et voluntáte labiórum ejus non fraudásti eum. \textred{℣.}~Posuísti Dómine super caput ejus corónam de lápide pretióso. \textred{E}t.

\subsubsection*{\center \normalfont \textred{Lectio v}}
\lettrine{H}abémus, delictíssimi, in vita Joánnis, et quod digne mirémur, et quod salúbrius imitémur. Quis non omni admiratióne stupéndum ducat, carne circúmdatum, cœlis paténtibus, carnis óculos supérno infígere lúmini, et vidére visiónes Dei? Hæc, atque simília in sancto hómine venerámur. Mérito quidem quando et sine salútis perículo possunt non fíeri, et sine salútis perículo néqueunt usurpári. Tútius sane æmulánda solidióra, quam sublimióra, et quæ magis virtútem redóleant, glóriam minus. Studeámus proínde móribus conformári.

\textred{℟.}~Amávit eum Dóminus, et ornávit eum~: stolam glóriæ índuit eum~: \textred{*}~Et ad portas paradísi coronávit eum. \textred{℣.}~Induit eum Dóminus lorícam fídei, et ornávit eum. \textred{E}t.

\subsubsection*{\center \normalfont \textred{Lectio vj}}
\lettrine{Æ}mulémur in viro sóbrium victum, devótum afféctum, æmulémur mansuetúdinem, spíritus, castimóniam córporis, oris custódiam, ánimi puritátem, pónere frænum iræ, et modum linguæ, dormíre párcius, oráre frequéntius, commonére nosmetípsos psalmis, hymnis, cánticis spirituálibus~: diébus júngere noctes, et divínis láudibus occupári. Æmulémur charísmata melióra: discámus ab ipso, quod mitis, et húmilis corde fuit. Æmulémur, inquam, quod éxstitit jucúndus ad hóspites, pátiens ad peccántes, benígnus ad omnes. Hoc enim mélius~: in his forma est, cui imprimámur~: in miráculis glória, a qua reprimámur~: Illa lætíficent, ista ædíficent; móveant illa, ista promóveant.

\textred{℟.}~Iste sanctus digne in memóriam vértitur hóminum~: qui ad gáudium tránsiit Angelórum~: quóniam in hac peregrinatióne solo córpore constitútus, cogitatióne, et aviditáte, \textred{*}~In illa ætérna pátria conversátus est. \textred{℣.}~Vínculis carnis absolútus, taléntum sibi créditum Dómino suo duplicátum reportávit. \textred{I}n illa ætérna. Glória Patri. \textred{I}n illa ætérna.

\subsubsection*{\center \textsc{\small IN III NOCTURNO}\\ 
\normalfont Léctio sancti Evangélii\\ secúndum Lucam\\
\textred{Lectio vij \hfill Cap.~11,~33-36}}
\lettrine{I}n illo témpore~: Dixit Jesus turbis~: Nemo lucérnam accéndit, et in abscóndito ponit, neque sub módio, sed supra candelábrum; ut, qui ingrediúntur lumen vídeant. Et réliqua.

\begin{center}Homilía sancti Joánnis Chrysóstomi\\
\small \textred{Hom. 10. in cap. 5. Matth.}\end{center}
\lettrineThree{P}{er} hanc comparatiónem vult osténdere, quare Sanctos suos ipse Christus maniféstat, et non sinit eos esse abscónsos, id est, quia nec qui accéndunt lucérnam, ídeo eam accéndunt, ut ponant sub módio, sed super candelábrum, ut lúceat ómnibus, qui in domo sunt. Qui sunt accensóres lucérnæ? Pater, et Fílius. Quæ est illa lucérna? Verbum divínum, de quo dictum est~: Lucérna pédibus meis verbum tuum. Quod est candelábrum? Ecclésia, quæ bájulat verbum vitæ; consequénter dicit et Paulus~: inter quos lucétis, sicut luminária in sǽculo, verbum vitæ continéntes. Consequénter et omnis ecclesiásticus vir habens verbum Dei, dícitur candelábrum.

\textred{℟.}~Magnificávit eum Dóminus in conspéctu regum~: \textred{*}~Et dedit illi corónam glóriæ. \textred{℣.}~Státuit illi testaméntum sempitérnum. \textred{E}t.

\subsubsection*{\center \normalfont \textred{Lectio viij}}
\lettrine{M}odii autem sunt hómines mundiáles, vácui a Deo, et ab ómnibus, quæ sunt Dei. Dicúntur autem módii, quóniam sicut módii désuper quidem vácui sunt, desúbtus autem pleni~: sic omnes mundi amatóres, et carnis, in spirituálibus quidem et divínis, quæ superióra recte esse dicúntur, vácui sunt, et insensáti~: in inferióribus autem, id est, in mundiálibus, et terrénis aliquátenus pleni vidéntur, et sapiéntes. Proptérea sicut in módio si volúeris áliquid repónere, necésse est, ut cadat deórsum ad inferióra ejus, et tenebrósa~: sic hómini carnáli, et mundiáli, et vácuo, si díxeris áliquid, quod est Dei, non illud tenet.

\textred{℟.}~Vir Israelíta, gaude cohéres Christi: \textred{*}~Ora pro nobis. \textred{℣.}~Ut précibus tuis a Deo véniam sperémus, devoti ad te confúgimus. \textred{O}ra.

\subsubsection*{\center \normalfont \textred{Lectio ix}}
\lettrine{N}unc autem sciens, quóniam ille est laudábilis, qui facit quod docet, addit, et dicit; sic lúceat lux vestra coram homínibus, ut vídeant ópera vestra bona, et magníficent Patrem vestrum, qui in cœlis est~: id est, sic illumináte, non ut verba vestra tantúmmodo áudiant hómines, sed ut ópera vestra bona vídeant, ut quos illuminavéritis per verbum, condiátis per exémpla óperum vestrórum~: quóniam qui docet, et facit quod docet, vere ille docet~: qui autem non facit quod docet, non álium docet, sed seípsum condémnat. Et mélius est fácere et non docére, quam docére et non fácere; quóniam qui facit, et si tacúerit, áliquos córrigit suo exémplo.

\textred{℟.}~Sancte Joánnes Christi Conféssor, audi rogántes sérvulos, \textred{*}~Et impetrátam nobis cœlitus tu defer indulgéntiam. \textred{℣.}~O sancte Joánnes sidus áureum Dómini grátia, servórum gémitus sólita súscipe cleméntia. \textred{E}t. Glória Patri. \textred{E}t.
 
\smallskip
Te Deum laudámus. %\textred{6}

\textred{Sacerd. ℣.}~Os justi meditábitur sapiéntiam. \textred{℟.}~Et lingua ejus loquétur judícium, allelúja.

\smallskip
{\small \textred{In II Vesperis, commemoratio S. Bibianæ:}}

\textred{Ant.} Accínxit * fortitúdine lumbos suos, et roborávit bráchium suum~: ideóque lucérna ejus non exstinguétur in sempitérnum.

\textred{℣.}~Diffúsa est grátia in lábiis tuis. \textred{℟.}~Proptérea benedíxit te Deus in ætérnum.

\smallskip
Orémus. \hfill \textred{Oratio}
\lettrine{D}eus ómnium largítor bonórum, qui in fámula tua Bibiána, cum virginitátis flore martýrii palmam cunjunxísti; mentes nostras, ejus intercessióne, tibi caritáte conjúnge, ut amótis perículis, prǽmia consequámur ætérna. Per Dóminum.

\begin{center}\textcolor{red}{\rule{\linewidth}{0.5pt}} \end{center}

\end{multicols}

\end{document}